%%%%%%%%%%%%%%%%
%%% num 16 %%%
%%%%%%%%%%%%%%%%

\Chapter{16}

\Verse{1}
{
	\hRJE{וַיִּקַּח קֹרַח בֶּן יִצְהָר בֶּן קְהָת בֶּן לֵוִי וְדָתָן וַאֲבִירָם בְּנֵי אֱלִיאָב וְאוֹן בֶּן פֶּלֶת בְּנֵי רְאוּבֵן׃}
}
{
	\eRJE{
		And Korah son of Izhar son of Kohath son of Levi, and
		Dathan and Abiram, sons of Eliab, and On, son of Peleth,
		sons of Reuben, took [image]
	}
}
{
%	\footnote{
%		Korah. He is a first cousin of Moses and Aaron (Exod 6:18–21). Footnote
%		16:1. Korah … took. He assembles a group with him, and he frames his
%		argument in terms of the people (“all of them are holy”). He is the
%		prototype of some politicians, religious leaders, and others, who act as
%		though they represent a group and are working for the people—when their
%		motives are suspect.
%	}
}

\Verse{2}
{
	\hRJE{וַיָּקֻמוּ לִפְנֵי מֹשֶׁה וַאֲנָשִׁים מִבְּנֵי יִשְׂרָאֵל חֲמִשִּׁים וּמָאתָיִם נְשִׂיאֵי עֵדָה קְרִאֵי מוֹעֵד אַנְשֵׁי שֵׁם׃}
}
{
	\eRJE{
		and got up in front of Moses—and two hundred fifty people
		from the children of Israel, chieftains of the
		congregation, prominent ones of the assembly, people of
		repute.
	}
}
{
}

\Verse{3}
{
	\hRJE{וַיִּקָּהֲלוּ עַל מֹשֶׁה וְעַל אַהֲרֹן וַיֹּאמְרוּ אֲלֵהֶם רַב לָכֶם כִּי כל הָעֵדָה כֻּלָּם קְדֹשִׁים וּבְתוֹכָם יְהֹוָה וּמַדּוּעַ תִּתְנַשְּׂאוּ עַל קְהַל יְהֹוָה׃}
}
{
	\eRJE{
		And they assembled against Moses and against Aaron and
		said to them, “You have much! Because all of the
		congregation, all of them, are holy, and YHWH is among
		them. And why do you raise yourselves up over YHWH’s
		community?”
	}
}
{
%	\footnote{
%		all of the congregation, all of them, are holy. It is true that all the
%		people have been described as a holy nation. Korah uses this now to
%		enhance his claim that he has as much right to the priesthood as Aaron.
%		But it is a deception, mixing two different meanings and degrees of
%		holiness. It is like telling a child that he or she is now “big” or
%		“grown up”—and then the child demands all the rights and privileges of
%		grownups.
%	}
}

\Verse{4}
{
	\hRJE{וַיִּשְׁמַע מֹשֶׁה וַיִּפֹּל עַל פָּנָיו׃}
}
{
	\eRJE{And Moses listened, and he fell on his face. [image]}
}
{
}

\Verse{5}
{
	\hRJE{וַיְדַבֵּר אֶל קֹרַח וְאֶל כּל עֲדָתוֹ לֵאמֹר בֹּקֶר וְיֹדַע יְהֹוָה אֶת אֲשֶׁר לוֹ וְאֶת הַקָּדוֹשׁ וְהִקְרִיב אֵלָיו וְאֵת אֲשֶׁר יִבְחַר בּוֹ יַקְרִיב אֵלָיו׃}
}
{
	\eRJE{
		And he spoke to Korah and to all of his congregation,
		saying, “In the morning YHWH will make known who is His
		and who is holy, and He will bring him close to Him. And
		He will bring the one He chooses close to Him.
	}
}
{
%	\footnote{
%		who is holy. Moses sets up a test of who is holy, to establish what it
%		means to be holy in the priestly sense, as opposed to the general
%		concept that the people should strive to be holy. The test will be the
%		use of incense, because only the holy—namely, priests—are permitted to
%		burn incense. The test will establish this, and it will be declared
%		formally following this episode (Num 17:5). And Korah and his company
%		cannot be unaware of the significance of the correct burning of incense,
%		because it was precisely a misuse of incense that led to the death of
%		Aaron’s sons, Nadab and Abihu (Lev 10:1–2).
%	}
}

\Verse{6}
{
	\hRJE{זֹאת עֲשׂוּ קְחוּ לָכֶם מַחְתּוֹת קֹרַח וְכל עֲדָתוֹ׃}
}
{
	\eRJE{
		Do this: Take incense burners, Korah and all his
		congregation,
	}
}
{
%	\footnote{
%		Do this. Why does Korah go along with the test? Do he and his followers
%		really believe that God would support them? Answer: they have no choice.
%		Moses does not propose a test. He commands Korah and his group. His
%		verbs are four imperatives in a row: do, take, put, set! And Korah
%		cannot refuse. The whole point, after all, is that Korah is seeking
%		priesthood (Num 16:10). What Moses challenges him to do is to perform
%		one of the acts that he would commonly have to do if he were to become a
%		priest: burn incense. Korah’s claim on the priesthood would evaporate if
%		he were to decline the opportunity he has been offered to do a priestly
%		task.
%	}
}

\Verse{7}
{
	\hRJE{וּתְנוּ בָהֵן אֵשׁ וְשִׂימוּ עֲלֵיהֶן קְטֹרֶת לִפְנֵי יְהֹוָה מָחָר וְהָיָה הָאִישׁ אֲשֶׁר יִבְחַר יְהֹוָה הוּא הַקָּדוֹשׁ רַב לָכֶם בְּנֵי לֵוִי׃}
}
{
	\eRJE{
		and put fire in them and set incense on them in front of
		YHWH tomorrow. And it will be that the man whom YHWH will
		choose, he will be the holy one. You have much, sons of
		Levi!”
	}
}
{
}

\Verse{8}
{
	\hRJE{וַיֹּאמֶר מֹשֶׁה אֶל קֹרַח שִׁמְעוּ נָא בְּנֵי לֵוִי׃}
}
{
	\eRJE{And Moses said to Korah, “Listen, sons of Levi, [image]}
}
{
}

\Verse{9}
{
	\hRJE{הַמְעַט מִכֶּם כִּי הִבְדִּיל אֱלֹהֵי יִשְׂרָאֵל אֶתְכֶם מֵעֲדַת יִשְׂרָאֵל לְהַקְרִיב אֶתְכֶם אֵלָיו לַעֲבֹד אֶת עֲבֹדַת מִשְׁכַּן יְהֹוָה וְלַעֲמֹד לִפְנֵי הָעֵדָה לְשָׁרְתָם׃}
}
{
	\eRJE{
		is it too small a thing for you that Israel’s {\God} has
		distinguished you from the congregation of Israel to bring
		you close to him, to do the work of YHWH’s Tabernacle and
		to stand in front of the congregation to minister for
		them,
	}
}
{
}

\Verse{10}
{
	\hRJE{וַיַּקְרֵב אֹתְךָ וְאֶת כּל אַחֶיךָ בְנֵי לֵוִי אִתָּךְ וּבִקַּשְׁתֶּם גַּם כְּהֻנָּה׃}
}
{
	\eRJE{
		and that He has brought you and all your brothers, the
		sons of Levi, with you close to Him? And you seek
		priesthood as well?!
	}
}
{
}

\Verse{11}
{
	\hRJE{לָכֵן אַתָּה וְכל עֲדָתְךָ הַנֹּעָדִים עַל יְהֹוָה וְאַהֲרֹן מַה הוּא כִּי (תלונו) [תַלִּינוּ] עָלָיו׃}
}
{
	\eRJE{
		Therefore you and all your congregation who are gathering
		are against YHWH! And Aaron, what is he that you complain
		against him?” [image]
	}
}
{
}

\Verse{12}
{
	\hRJE{וַיִּשְׁלַח מֹשֶׁה לִקְרֹא לְדָתָן וְלַאֲבִירָם בְּנֵי אֱלִיאָב וַיֹּאמְרוּ לֹא נַעֲלֶה׃}
}
{
	\eRJE{
		And Moses sent to call Dathan and Abiram, sons of Eliab,
		and they said, “We won’t come up.
	}
}
{
%	\footnote{
%		Moses sent to call Dathan and Abiram. Two groups with two different
%		complaints have joined together to challenge Moses. Korah’s company
%		challenge Aaron’s hold on the priesthood; Dathan and Abiram challenge
%		Moses’ hold on the leadership of Israel. This political union of two
%		agendas is presumably arranged to strengthen the positions of both.
%		Moses, however, keeps them separate. He has responded first to Korah and
%		his company who are challenging the priestly prerogatives of Moses and
%		Aaron. Now he turns to Dathan and Abiram and summons them. By separating
%		them, he weakens their position.
%	}
}

\Verse{13}
{
	\hRJE{הַמְעַט כִּי הֶעֱלִיתָנוּ מֵאֶרֶץ זָבַת חָלָב וּדְבַשׁ לַהֲמִיתֵנוּ בַּמִּדְבָּר כִּי תִשְׂתָּרֵר עָלֵינוּ גַּם הִשְׂתָּרֵר׃}
}
{
	\eRJE{
		Is it a small thing that you brought us up from a land
		flowing with milk and honey to kill us in the wilderness,
		that you lord it over us as well?
	}
}
{
%	\footnote{
%		you brought us up from a land flowing with milk and honey. They are
%		referring to Egypt! They use the rhetorical device of reversing their
%		opponent’s position: they say that Moses has taken them out of a land
%		flowing with milk and honey. It is a clever device (common in political
%		debate). And it is outrageous: accusing the liberator of being the
%		source of their problems, picturing the land of slavery as the land of
%		milk and honey, and, as usual, focusing on Moses rather than on God. And
%		then they say, “Will you put out those people’s eyes?”—accusing Moses of
%		being the one who is deceiving the people, when it is they who are the
%		deceivers. This is a battle for leadership, which is to say it is a
%		political battle, which is to say it is a battle for power; and it
%		dramatizes the danger of the power-seeking politician who is a skillful
%		speaker. Moses defeats his opponents because God intervenes with a
%		miracle. We who cannot count on earthquakes to identify the corrupt ones
%		among us must be thoughtful, sensitive, and diligent when we listen to
%		political argument.
%	}
}

\Verse{14}
{
	\hRJE{אַף לֹא אֶל אֶרֶץ זָבַת חָלָב וּדְבַשׁ הֲבִיאֹתָנוּ וַתִּתֶּן לָנוּ נַחֲלַת שָׂדֶה וָכָרֶם הַעֵינֵי הָאֲנָשִׁים הָהֵם תְּנַקֵּר לֹא נַעֲלֶה׃}
}
{
	\eRJE{
		Besides, you haven’t brought us to a land flowing with
		milk and honey or given us possession of field or
		vineyard. Will you put out those people’s eyes? We won’t
		come up.”
	}
}
{
%	\footnote{
%		you haven’t brought us to a land flowing with milk and honey or given us
%		possession of field or vineyard. This is true. Part of a persuasive lie
%		is to mix in something that is true.
%	}
}

\Verse{15}
{
	\hRJE{וַיִּחַר לְמֹשֶׁה מְאֹד וַיֹּאמֶר אֶל יְהֹוָה אַל תֵּפֶן אֶל מִנְחָתָם לֹא חֲמוֹר אֶחָד מֵהֶם נָשָׂאתִי וְלֹא הֲרֵעֹתִי אֶת אַחַד מֵהֶם׃}
}
{
	\eRJE{
		And Moses was very angry, and he said to YHWH, “Don’t turn
		to their offering. Not one ass of theirs have I taken
		away, and I haven’t wronged one of them.”
	}
}
{
}

\Verse{16}
{
	\hRJE{וַיֹּאמֶר מֹשֶׁה אֶל קֹרַח אַתָּה וְכל עֲדָתְךָ הֱיוּ לִפְנֵי יְהֹוָה אַתָּה וָהֵם וְאַהֲרֹן מָחָר׃}
}
{
	\eRJE{
		And Moses said to Korah, “You and all your congregation,
		be in front of YHWH—you and they and Aaron—tomorrow.
		{[image]}
	}
}
{
}

\Verse{17}
{
	\hRJE{וּקְחוּ אִישׁ מַחְתָּתוֹ וּנְתַתֶּם עֲלֵיהֶם קְטֹרֶת וְהִקְרַבְתֶּם לִפְנֵי יְהֹוָה אִישׁ מַחְתָּתוֹ חֲמִשִּׁים וּמָאתַיִם מַחְתֹּת וְאַתָּה וְאַהֲרֹן אִישׁ מַחְתָּתוֹ׃}
}
{
	\eRJE{
		And each man take his fire-holder, and put incense on
		them, and each man bring his fire-holder forward in front
		of YHWH, two hundred fifty fire-holders, and you and
		Aaron, each man his fire-holder.”
	}
}
{
}

\Verse{18}
{
	\hRJE{וַיִּקְחוּ אִישׁ מַחְתָּתוֹ וַיִּתְּנוּ עֲלֵיהֶם אֵשׁ וַיָּשִׂימוּ עֲלֵיהֶם קְטֹרֶת וַיַּעַמְדוּ פֶּתַח אֹהֶל מוֹעֵד וּמֹשֶׁה וְאַהֲרֹן׃}
}
{
	\eRJE{
		And each man took his fire-holder, and they put fire on
		them and set incense on them, and they stood at the
		entrance of the Tent of Meeting, and Moses and Aaron.
		{[image]}
	}
}
{
}

\Verse{19}
{
	\hRJE{וַיַּקְהֵל עֲלֵיהֶם קֹרַח אֶת כּל הָעֵדָה אֶל פֶּתַח אֹהֶל מוֹעֵד וַיֵּרָא כְבוֹד יְהֹוָה אֶל כּל הָעֵדָה׃}
}
{
	\eRJE{
		And Korah assembled all the congregation against them, to
		the entrance of the Tent of Meeting. And YHWH’s glory
		appeared to all the congregation.
	}
}
{
}

\Verse{20}
{
	\hRJE{וַיְדַבֵּר יְהֹוָה אֶל מֹשֶׁה וְאֶל אַהֲרֹן לֵאמֹר׃}
}
{
	\eRJE{And YHWH spoke to Moses and to Aaron, saying, [image]}
}
{
}

\Verse{21}
{
	\hRJE{הִבָּדְלוּ מִתּוֹךְ הָעֵדָה הַזֹּאת וַאֲכַלֶּה אֹתָם כְּרָגַע׃}
}
{
	\eRJE{
		“Separate from among this congregation, and I’ll finish
		them in an instant!”
	}
}
{
}

\Verse{22}
{
	\hRJE{וַיִּפְּלוּ עַל פְּנֵיהֶם וַיֹּאמְרוּ אֵל אֱלֹהֵי הָרוּחֹת לְכל בָּשָׂר הָאִישׁ אֶחָד יֶחֱטָא וְעַל כּל הָעֵדָה תִּקְצֹף׃}
}
{
	\eRJE{
		And they fell on their faces and said, “{\God}, the {\God} of
		the spirits of all flesh, will one man sin and you be
		angry at all the congregation?”
	}
}
{
%	\footnote{
%		will one man sin and you be angry at all the congregation. Korah has
%		brought the whole congregation of Israel along (16:19), and God has
%		threatened to destroy all of them (16:21), and so Moses and Aaron plead
%		that only those who have done something wrong should suffer. “One man”
%		here cannot mean literally only a single person; it must refer to every
%		individual in Korah’s congregation, as opposed to the entire
%		congregation of Israel. It is reminiscent of Abraham’s plea in the case
%		of Sodom and Gomorrah (“Will you also annihilate the virtuous with the
%		wicked?” Gen 18:23). The connection to the case of Sodom is highlighted
%		by the fact that Moses then tells the people to get away “or else you’ll
%		be annihilated,” which are the words that the angels say at Sodom—and
%		these words occur nowhere else in the Torah (see Gen 18:23,24;
%		19:15,17).
%	}
}

\Verse{23}
{
	\hRJE{וַיְדַבֵּר יְהֹוָה אֶל מֹשֶׁה לֵּאמֹר׃}
}
{
	\eRJE{And YHWH spoke to Moses, saying, [image]}
}
{
}

\Verse{24}
{
	\hRJE{דַּבֵּר אֶל הָעֵדָה לֵאמֹר הֵעָלוּ מִסָּבִיב לְמִשְׁכַּן קֹרַח דָּתָן וַאֲבִירָם׃}
}
{
	\eRJE{
		“Speak to the congregation, saying, ‘Get up from around
		the tabernacle of Korah, Dathan and Abiram.’”
	}
}
{
}

\Verse{25}
{
	\hRJE{וַיָּקם מֹשֶׁה וַיֵּלֶךְ אֶל דָּתָן וַאֲבִירָם וַיֵּלְכוּ אַחֲרָיו זִקְנֵי יִשְׂרָאֵל׃}
}
{
	\eRJE{
		And Moses got up and went to Dathan and Abiram, and
		Israel’s elders went after him.
	}
}
{
}

\Verse{26}
{
	\hRJE{וַיְדַבֵּר אֶל הָעֵדָה לֵאמֹר סוּרוּ נָא מֵעַל אהֳלֵי הָאֲנָשִׁים הָרְשָׁעִים הָאֵלֶּה וְאַל תִּגְּעוּ בְּכל אֲשֶׁר לָהֶם פֶּן תִּסָּפוּ בְּכל חַטֹּאתָם׃}
}
{
	\eRJE{
		And he spoke to the congregation, saying, “Turn away from
		the tents of these wicked men and don’t touch anything
		that is theirs, or else you’ll be annihilated through all
		their sins.”
	}
}
{
}

\Verse{27}
{
	\hRJE{וַיֵּעָלוּ מֵעַל מִשְׁכַּן קֹרַח דָּתָן וַאֲבִירָם מִסָּבִיב וְדָתָן וַאֲבִירָם יָצְאוּ נִצָּבִים פֶּתַח אהֳלֵיהֶם וּנְשֵׁיהֶם וּבְנֵיהֶם וְטַפָּם׃}
}
{
	\eRJE{
		And they got up from around the tabernacle of Korah,
		Dathan and Abiram. And Dathan and Abiram came out,
		standing at the entrance of their tents, and their wives
		and their children and their infants.
	}
}
{
}

\Verse{28}
{
	\hRJE{וַיֹּאמֶר מֹשֶׁה בְּזֹאת תֵּדְעוּן כִּי יְהֹוָה שְׁלָחַנִי לַעֲשׂוֹת אֵת כּל הַמַּעֲשִׂים הָאֵלֶּה כִּי לֹא מִלִּבִּי׃}
}
{
	\eRJE{
		And Moses said, “By this you’ll know that YHWH sent me to
		do all these things, because it’s not from my own heart.
		{[image]}
	}
}
{
}

\Verse{29}
{
	\hRJE{אִם כְּמוֹת כּל הָאָדָם יְמֻתוּן אֵלֶּה וּפְקֻדַּת כּל הָאָדָם יִפָּקֵד עֲלֵיהֶם לֹא יְהֹוָה שְׁלָחָנִי׃}
}
{
	\eRJE{
		If these die like the death of every human, and the event
		of every human happens to them, then YHWH hasn’t sent me.
		{[image]}
	}
}
{
}

\Verse{30}
{
	\hRJE{וְאִם בְּרִיאָה יִבְרָא יְהֹוָה וּפָצְתָה הָאֲדָמָה אֶת פִּיהָ וּבָלְעָה אֹתָם וְאֶת כּל אֲשֶׁר לָהֶם וְיָרְדוּ חַיִּים שְׁאֹלָה וִידַעְתֶּם כִּי נִאֲצוּ הָאֲנָשִׁים הָאֵלֶּה אֶת יְהֹוָה׃}
}
{
	\eRJE{
		But if YHWH will create something, and the ground will
		open its mouth and swallow them and all that they have,
		and they’ll go down alive to Sheol, then you’ll know that
		these people have rejected YHWH.”
	}
}
{
}

\Verse{31}
{
	\hRJE{וַיְהִי כְּכַלֹּתוֹ לְדַבֵּר אֵת כּל הַדְּבָרִים הָאֵלֶּה וַתִּבָּקַע הָאֲדָמָה אֲשֶׁר תַּחְתֵּיהֶם׃}
}
{
	\eRJE{
		And it was as he was finishing speaking all these things,
		and the ground that was under them was broken up,
	}
}
{
}

\Verse{32}
{
	\hRJE{וַתִּפְתַּח הָאָרֶץ אֶת פִּיהָ וַתִּבְלַע אֹתָם וְאֶת בָּתֵּיהֶם וְאֵת כּל הָאָדָם אֲשֶׁר לְקֹרַח וְאֵת כּל הָרְכוּשׁ׃}
}
{
	\eRJE{
		and the earth opened its mouth and swallowed them and
		their households and all the people who were with Korah
		and all the property,
	}
}
{
%	\footnote{
%		them and their households. Meaning: their entire families die with them.
%		Why does the entire household suffer? Elsewhere in the law we are told
%		that each person is responsible for his or her own crimes. Children are
%		not to suffer for their parents’ acts (Deut 24:16). But that principle
%		applies only to ethical law: crimes between one human and another. In
%		ritual law—violations against God, such as blasphemy or, in this case,
%		openly opposing an instruction from God (“these people have rejected
%		YHWH”)—anyone who comes in contact with the source of the offense is
%		tainted by it, regardless of that person’s own intent. Thus, for
%		example, in the book of Joshua, Achan keeps some of the forbidden spoils
%		of Jericho, which are designated for complete destruction ([image]erem).
%		He conceals them in his tent. When his offense is revealed he and the
%		spoils “and his sons and his daughters and his ox and his ass and his
%		flock and his tent and everything he had” are stoned and burned (Josh
%		7:24). It is difficult to understand in the postbiblical age, but the
%		concepts of ritual, sacred realms, purity and impurity, and violation of
%		sacred zones are fundamental to Israel’s life as pictured in the Bible’s
%		early books. It may be that the advance in psychology and especially of
%		psychoanalysis in this century can help us to appreciate the mind-set of
%		those who conceived of this notion of guilt by contact. We can recognize
%		that people feel guilt even about things that are not their fault
%		morally. Victims of childhood incest frequently feel guilt through their
%		adult lives. As a closer example: as a child I observed another boy tear
%		a Torah scroll, and I recall how upsetting this was to him and others
%		even though he did not necessarily do it out of lack of care. I also
%		recall, from my student days, knocking a pious rabbi’s black hat from a
%		hook onto the floor and feeling terrible even though it was an accident.
%		(He was wise; he just said lightly, “Praise God that my head wasn’t in
%		it at the time.”) The point is that we do have feelings of guilt in
%		connection with things that we hold in awe. And this is multiplied a
%		thousandfold in the narrative of the age of the Torah: when the divine
%		is believed to be close at hand.
%	}
}

\Verse{33}
{
	\hRJE{וַיֵּרְדוּ הֵם וְכל אֲשֶׁר לָהֶם חַיִּים שְׁאֹלָה וַתְּכַס עֲלֵיהֶם הָאָרֶץ וַיֹּאבְדוּ מִתּוֹךְ הַקָּהָל׃}
}
{
	\eRJE{
		and they went down, they and all that they had, alive to
		Sheol. And the earth covered them over, and they perished
		from among the community.
	}
}
{
}

\Verse{34}
{
	\hRJE{וְכל יִשְׂרָאֵל אֲשֶׁר סְבִיבֹתֵיהֶם נָסוּ לְקֹלָם כִּי אָמְרוּ פֶּן תִּבְלָעֵנוּ הָאָרֶץ׃}
}
{
	\eRJE{
		And all Israel that was around them fled at the sound of
		them, for they said, “Or else the earth will swallow us.”
		{[image]}
	}
}
{
}

\Verse{35}
{
	\hRJE{וְאֵשׁ יָצְאָה מֵאֵת יְהֹוָה וַתֹּאכַל אֵת הַחֲמִשִּׁים וּמָאתַיִם אִישׁ מַקְרִיבֵי הַקְּטֹרֶת׃}
}
{
	\eRJE{
		And fire had gone out from YHWH and consumed the two
		hundred fifty people offering the incense.
	}
}
{
%	\footnote{
%		fire had gone out. This is worded in the past perfect. (The noun
%		precedes the verb in the Hebrew.) That is, this must be understood to
%		have happened to Korah’s incense-burning group before the earthquake
%		swallowed them along with Dathan’s and Abiram’s households.
%	}
}
